
\section{\textbf{\texttt{search()}} ORM Method}

\begin{frame}[fragile]{\textbf{\texttt{search()}} Method}
	\begin{itemize}
		\item The \textbf{\texttt{search()}} method is an ORM method used to \textbf{find records} that match a domain (filter condition).
		\item It is one of the most used ORM methods in Odoo.
		\item The method returns a recordset of matching records.
\begin{block}{Method Syntax}
	\begin{itemize}
		\item The standard syntax of the search() method is:
\begin{lstlisting}
@api.model
def search(self, domain, offset=0, limit=None, order=None):
	return super().search(domain, offset=offset, limit=limit, order=order)
\end{lstlisting}
		\item Important parts: \textbf{\texttt{domain}}, \textbf{\texttt{offset}}, \textbf{\texttt{limit}}, \textbf{\texttt{order}}.
	\end{itemize}
\end{block}
	\end{itemize}
\end{frame}

\begin{frame}[fragile]{Breaking Down the Method Signature}
	\begin{block}{}
		\begin{itemize}
			\item \textbf{\texttt{@api.model}}: Search runs at model level (not on one existing record).
			\item \textbf{\texttt{self}}: Represents the model, e.g.\ \texttt{student.student()}.
			\item \textbf{\texttt{domain}}: A list of conditions used to filter records.
\begin{lstlisting}
domain = [('age', '>=', 18), ('gender', '=', 'male')]
\end{lstlisting}
			\item Each condition is a tuple: \texttt{(field\_name, operator, value)}.
			\item Multiple conditions are combined with AND by default.
		\end{itemize}
	\end{block}
\end{frame}

\begin{frame}[fragile]{Breaking Down Domain Operators}
	Common operators used inside a domain:
	\begin{itemize}
		\item \texttt{=} equal
		\item \texttt{!=} not equal
		\item \texttt{>}, \texttt{>=}, \texttt{<}, \texttt{<=} comparisons
		\item \texttt{like} / \texttt{ilike} text search (\texttt{ilike} is case-insensitive)
		\item \texttt{in} / \texttt{not in} membership
		\item \texttt{child\_of} hierarchical search
	\end{itemize}
\begin{lstlisting}
[('name', 'ilike', 'ahmed')]
[('id', 'in', [15, 16, 17])]
['|', ('gender', '=', 'male'), ('age', '<', 18)]
\end{lstlisting}
	\begin{itemize}
		\item \texttt{'|'} means OR, \texttt{'\&'} means AND, \texttt{'!'} means NOT.
	\end{itemize}
\end{frame}

\begin{frame}[fragile]{Offset, Limit, and Order}
	\begin{itemize}
		\item \textbf{\texttt{offset}}: Skip the first N records (used for pagination).
		\item \textbf{\texttt{limit}}: Return at most N records.
		\item \textbf{\texttt{order}}: SQL-like ordering string.
\begin{lstlisting}
students = self.env['student.student'].search(
	[('active', '=', True)],
	offset=10,
	limit=5,
	order='name asc, id desc',
)
\end{lstlisting}
		\item This returns active students 11--15 when sorted by name.
	\end{itemize}
\end{frame}

\begin{frame}[fragile]{What Does \texttt{search()} Return?}
	\begin{block}{}
		\begin{itemize}
			\item The \textbf{\texttt{search()}} method returns a recordset.
\begin{lstlisting}
students = self.env['student.student'].search([
	('name', 'ilike', 'ahmed'),
])
print(students)
\end{lstlisting}
			Output is (example):
\begin{lstlisting}
student.student(15, 22, 31)
\end{lstlisting}
			\item If nothing matches, it returns an empty recordset: \texttt{student.student()}.
			\item An empty recordset is falsy in an \texttt{if} condition.
		\end{itemize}
	\end{block}
\end{frame}

\begin{frame}[fragile]{Method Call and Practical Example}
\begin{block}{Method Call}
\begin{lstlisting}
students = self.env['student.student'].search([
	('age', '>=', 18),
	('active', '=', True),
], limit=20, order='age desc')
\end{lstlisting}
	\begin{itemize}
		\item Here, \textbf{\texttt{.search(...)}} is the \textbf{method call}.
		\item Always prefer domains over building raw SQL for normal queries.
		\item Record rules and access rights are applied automatically.
	\end{itemize}
\end{block}
\end{frame}
